%% APPENDIX D -- the false-alarm expectation, the control ensemble, and the
%% execution-block quality criterion.  Carries \label{app:falsealarm}
%% (04_method.tex, 05_results.tex, 08_backmatter.tex) and
%% \label{sec:blockquality} (04_method.tex, 05_results.tex).
%%
%% ★★★★ WHAT THIS APPENDIX IS FOR, AND WHAT IT MAY NOT GROW BACK INTO.  Four
%% things and no more: the expectation's definition, the resampling the count
%% is compared through, the principal numbers, and the control model's
%% limitations.  Everything else that stood here is either stated once in
%% \S\ref{sec:trigunif} -- the extraction geometry, the two probabilities and
%% their correction, the bound's qualifications -- or is a diagnostic that
%% decided nothing: the threshold scan's margins and its counts at the next
%% and median blocks, the unscaled form of the block rule, the cell-grain
%% reading of the same expectation, the clean coarse windows' addition to it,
%% the systemic-velocity sensitivity, the near-band ratio read half-level by
%% half-level and its value at the trigger itself, the reserved blocks that
%% share no session, and the scrambling control.  Each was measured and each
%% is in the generators; none of them changes a disposition, and a limitation
%% is condensed here, never deleted.
\section{The false-alarm expectation and the control ensemble}
\label{app:falsealarm}

This appendix defines the chance expectation for the unattributed crossings,
gives the resampling the observed count is compared through, and states what
the control ensemble that calibrates it can and cannot support.

\subsection{The expectation, and what it is conditional on}

\emph{One attribution set serves both sides of the comparison}, the set of
\S\ref{sec:technosearch}: the mask of Appendix~\ref{app:linemask} within
$\pm\MaskHalfKms$\,km\,s$^{-1}$ in the star's own rest frame, which holds no
sulphur monoxide, so the \SoNCrossWord{} crossings that fall on SO
transitions are unattributed crossings and a control maximum on one of those
transitions is an unattributed exceedance.

Each of the \CeNWin{} Class~A windows of the primary census contributes two
quantities measured on itself. The fraction of its own \NCtrl{} control
positions that reaches the trigger is the probability that a position in that
field does so -- no cell count enters it and no form is assumed for the tail
-- and summed over the windows \CeExpTrig{} are expected to trigger by
chance. How many of those would be attributed is measured on the channels the
control maxima land on: \CeNExcAttr{} of \CeNExc{} exceedances fall inside
the mask, carried to the same star's rest frame through the same two steps
and against the same transition list as a cell at the stellar position, so
the attribution takes \CeMaskPct{} per cent of the expectation where a
maximum free to land anywhere on its own grid would take \CeUnifPct{} per
cent: the \CeNFullMask{} line-tuned windows lie wholly inside the mask, and a
uniform landing would under-attribute them. The atmospheric screen of
Appendix~\ref{app:rfi} is then applied to these same positions, on the same
radius and in the same topocentric frame as to the crossings, and takes
\CeTlCtrlCut{} of the \CeTlCtrlSurv{} exceedances the mask leaves. That
leaves \CeExp{} expected unattributed Class~A crossings against \CeObs{}
observed, as a sum over rows with nothing carried at an assumed rate. Both
sides are read over the whole primary beam \S\ref{sec:statistic} defines,
with each block on its own antenna geometry; \S\ref{sec:trigunif} reports
what a narrower field does to both of them, where the bound below becomes
\FrNarrowAddLim. System by system, each system's expectation formed from its
own control positions, \MtNSysLowP{} of the \MtNSysChance{} systems fall
below one chance in twenty against \MtNSysLowPExp{} expected. A further
\CeNWinTail{} searched Class~A windows lie outside the comparison because the
deposit does not tabulate them, together with the \CeNTailCrossWord{}
unattributed crossings toward \CeTailStar{} they hold, and enter neither
side.

\emph{A count is compared with the distribution of counts chance delivers,
and not with a range on its own mean.} A rate is drawn from \CeNBoot{}
resamplings -- each window's rate redrawn from its own \NCtrl{} probes, the
attributed share from its own \CeNExc{} exceedances, and the execution blocks
resampled because windows share them and are not independent trials -- and
the count scatters about that rate. A count of \CeObs{} or more then arises
with probability \CeProbNull{} at the control positions alone, a probability
conditional on the control-position model set out here and not on a null
distribution established independently of it. The same resampling bounds what
the comparison could have missed: added to each resampled rate, \CeAddLim{}
further crossings still leave a count as low as \CeObs{} credible at
\CeProbAtLim, and one more does not, at \CeProbAboveLim{} against the
\CeAlphaPct{} per cent level. So \CeAddLim{} is the largest added population
this comparison could not have rejected against the control positions' own
expectation, the uncertainty in \CeExp{} carried inside it by the same
resampling; it counts crossings and not transmitters, and several crossings
could come from one source.

\emph{The one assumption the searched data cannot test is that a star behaves
like a control position in its own field.} In the census the star carries
\CeNear{} times as many above-threshold cells as a control position of the
same field, measured in the band immediately below the trigger,
$\CeNearBandLo\le T<\CeNearBandHi$, with \CeNearIntLo--\CeNearIntHi{} allowed
at \CeCovPct{} per cent when the execution blocks are resampled; no cell
under test enters it, a band count being the difference of two cumulative
counts. Calibrating the stellar rate against that ratio rather than against
the control positions' own, as a sensitivity test, raises the expectation to
\CeExpNear{} and the probability to \CeProbNear{} and lowers the bound to
\CeAddLimNear: the conclusion that there is no significant excess is
unchanged under it, and the control positions' own rate is adopted as the
smaller of the two and the more demanding of the count under test. The
requirement runs the other way in any case -- the \CeObs{} observed would
stand above chance at the \CeAlphaPct{} per cent level only if stars reached
the trigger at \CeRCrit{} times a control position's rate or less, and every
ratio measured here is an excess instead. The reserved blocks measure the
same quantity where the search did not reach: their stars trigger in
\StHoObs{} of \StHoNFine{} fine-channel windows against \StHoExpFine{}
predicted, a ratio of \StHoRatio{} -- though every one of those events lies
in a block that shares an observing session with a searched one, and the
\HsIndBlocks{} reserved blocks that share none hold \HsIndObsWord{} trigger
\HsIndObsNoun{} between them, so they neither confirm the ratio nor
contradict it (Appendix~\ref{sec:heldout}).

Near a band edge the exceedance rate is $\times\VnEdgeRatio$ the interior
rate, pooled -- and the excess is one object: the window along the sightline
to HD~48370 contributes \VnEdgeDomPct{} per cent of every near-edge
exceedance in the survey, without it the ratio is
$\times\VnEdgeNoDomRatio$, and on a held-out half it is
$\times\VnEdgeHalfCleanRatio$ ($\pv{\VnEdgeHalfCleanP}$). No edge
correction is therefore applied to any limit.

\subsection{Execution-block quality}
\label{sec:blockquality}

The comparison assumes each block's control ensemble behaves like noise, and
one block makes that false: every one of its \NCtrl{} probes exceeds the
trigger, the smallest at \BqFailCtrlMin. A criterion is therefore applied to
every block and not to the one that prompted it, built from the median of a
window's own control ensemble, which separates a field with a source in it
from a field that is no null at all: a block is judged on that median divided
by the median over the windows of its own class, since a maximum over a
larger channel$\times$drift grid sits higher, and the same statistic flags
\DqNWin{} of the \NWindows{} windows individually at \BqFactor{} the survey
median, \DqThresh{} against \DqSurvMed. \emph{Neither form was fixed in
advance}: the condition was noticed in one block and the criterion written
afterwards, so it is post hoc, and what makes it trustworthy is how little
the position of the line matters -- the failing block sits at \BqFailQ{} and
the next most extreme at \BqNextQ, a factor \QsRatio{} apart, and over
\QsNStep{} thresholds spanning that interval the same single block is
selected every time. Of the \BqNBlocks{} primary-census blocks, \BqNFail{}
fails; its \BqFailNWin{} windows are all coarse-channel and hold the survey's
entire Class~B crossing population, \EvNCrossFlag{} of the \EvNCrossRaw{}
threshold crossings, which appear in the crossing ledger with their block
named and enter no count elsewhere.

\subsection{What the ensemble can resolve, and in what unit}

Ranking the star against \NCtrl{} control positions of the same field decides
which windows deserve examination, but $1/(\NCtrl+1)=\RankFloorVal$ is the
\emph{resolution} of the ensemble -- the smallest rank it can express -- and
not the probability that noise produced the window. A survey of \NWindows{}
windows would need a rank \BonfRankRatio{} times finer, and correlation makes
it coarser still: controls inside the synthesised beam are not independent,
so an annulus resolves no better than \RankEffLo--\RankEffHi{} however many
probes are drawn into it. That floor is the design's central limitation: no
single window here can be significant on the rank alone, which is why
confirmation requires independent recurrence. The rank distribution is weak
evidence in either direction for the same reason, a shift smaller than about
\VerKsPowerAca{} being invisible in a sample of this size.

The searched cells counted in \S\ref{sec:summary} are not independent trials
either: channel smoothing and a drift grid stepped twice per channel
correlate neighbouring ones, which alone would leave \StGridFrac{} of them
independent, while fitting each window's effective count against its own
control ensemble returns \StNindFrac{} and those ensembles' maxima sit at
\CtrlMaxRatioMed{} of the prediction for the maximum of the full searched
count treated as independent and Gaussian. Nothing turns on which is right,
because the expectation above is summed from measured rates and not from a
cell count: by the cell-count route the number reaching the trigger is
\StExpModel{} against \CeExpTrig, \StModelAgreePct{} per cent away. A
Gaussian tail over every searched cell predicts \CellExpect{} stellar
exceedances against the \NCross{} found; the attribution, the failing block,
the archival tail and the atmospheric screen account for the difference.

The widened recurrence window of \S\ref{sec:etacrv} is counted in the unit it
is measured in. A window covering more than one crossing is one
window and one measurement, since the statistic read is that window's own
maximum and does not depend on which crossing it is referred to: the
\WuNRead{} pairings of a crossing with a covering window are \WuNWin{}
distinct windows, \WuNMulti{} covering more than one. Of those, \WuNWinExc{}
hold a statistic above the trigger against \WuExpWin{} expected at
$P=\WuProbWin$, or \WuExpWinCtrl{} at $P=\WuProbWinCtrl$ when each window's
own control positions supply the rate in place of its cell count; the larger
probability is quoted. \WuNWinExcTwo{} are toward \WuTwoStars{} alone, and
over the other \WuNWinRest{} windows the comparison is \WuNWinRestExc{}
against \WuExpRest{} at $P=\WuProbRest$, so what departure there is is
concentrated and not general.

%% The chance expectation here is summed from the measured control exceedance
%% rates and each window's own masked grid share (chance_v414.py, round 310),
%% so no cell count and no Gaussian tail enters it and nothing in it is
%% conditioned on the crossings; the cell-count route is a cross-check only,
%% and the rank-conditioned expectation (\ClustStageMean against
%% \ClustStageNObs) is a pipeline diagnostic that must not be reachable as
%% "the" expectation.  tab:ctrlsep, tab:excess, tab:arraysplit and
%% fig:ctrldiag are generated and nothing cites them; r15-deposit, 2026-10-08:
%% they used to be described as "in the data release", and there is no data
%% release, so they are simply unpublished diagnostics.
\subsection{The registration of the predicted cell}
\label{sec:regbudget}

The recurrence test of \S\ref{sec:etacrv} reads one channel at one drift
rate, and the cell is registered between two epochs to \RgBudgetChMed{} of a
channel in every window -- half a channel for the prediction, half for the
discovery channel it is referred to, and smaller terms for the drift, the
systemic velocity and the ephemeris -- while the frame correction, the one
term that is large in channels rather than small, moves it by up to
\ZtFrameChanMax{} channels, and the stellar-frame centroid of a real
recurring line sits a median \RgCenMedKms\,km\,s$^{-1}$ from the predicted
cell over \RgCenNBlock{} block readings of the line-attributed recoveries.
Each reading was therefore re-formed from the worst cell within
\RgTolKms\,km\,s$^{-1}$ of the predicted one with the threshold raised to
\RgThrMax{} so that the wider reading keeps the single cell's own false-alarm
probability, it was pinned at every anchor across the span the discovery track
covers, which moves it by at most \AnSpanMaxKms\,km\,s$^{-1}$, and the bar
itself was reset to \ZtTrialThr{} to hold the chance of one false confirmation
across all \RcNRepeatMeas{} readings at \ZtFwPct{} per cent -- and no
disposition rests on any of the three, each crossing's own tolerance, cell
count, threshold and resulting bound accompanying this paper as data. What the
recoveries establish about the registration of a single cell rests on the
\RgRegNUnres{} readings of \RgRegStar{} alone, whose line is confined to one
channel; the rest test the frame correction only to the width of a resolved
line.

\subsection{From the widened window's exceedances to the pairs}
\label{sec:pairacct}

Two counts of \S\ref{sec:etacrv} are counts of different objects. The
\WdNExc{} above-trigger readings are window readings, one for each pairing of
an unattributed crossing with a window that covers its frequency in another
block -- \WdNCov{} such pairings over the \PaNCovWinCross{} crossings a later
block covers -- and they fall in \PaNRepWin{} windows of \PaNRepBlock{}
blocks, so a pair whose two members are both unattributed crossings is read
twice: \PaNPairTwice{} of the readings are such a second reading,
representing \PaNPairRead{} distinct pairs, of which \PaNPairBeyond{} lie
outside the interval the pair test admits -- the median covering window's
half-width, by up to \PaDvBeyondMax\,km\,s$^{-1}$ -- and \PaNPairInside{}
remain, which is the set Table~\ref{tab:pairs} is formed over;
Table~\ref{tab:pairacct} gives the accounting crossing by crossing. The observatory's own line-of-sight acceleration is removed from both
fitted drifts before they are compared, at most \TjFrameMaxHzS\,Hz\,s$^{-1}$
here. Of the \TjNPower{} pairs whose rates the grid resolves, every one fails
a drift changing at a constant rate, by \TjSigMin{} standard deviations in
the mildest case; the offset
$\frac{1}{2}(\dot\nu_1+\dot\nu_2)\Delta t$ such a
trajectory predicts has no freedom left once both rates and the interval are
fixed, and the one pair the searched grid itself separates is \EcGridStar's,
whose crossings lie \EcGridDvKms\,km\,s$^{-1}$ apart in blocks
\EcGridDtH\,hours apart, a single carrier having had to drift at
\EcGridFactor{} times the grid's ceiling between them. Nothing measured here
bounds the curvature in between, the ceiling bounding the line-of-sight
component of a carrier's acceleration and not how fast it can turn over;
the velocity excursion \emph{is} bounded, to within \EcVsupFac{} times the
circular amplitude over every eccentricity and orientation, which is why the
widened window of \S\ref{sec:etacrv} can be set against that excursion and a
reversal time cannot. The control row is a line that does not drift: the
\TjCtlLine{} line of \TjCtlStar{} is recovered at one stellar-frame frequency
\TjCtlDtD{} days apart and \TjCtlNSig{} standard deviations from the
prediction, so the comparison resolves a mean drift to
\TjCtlDriftResHzS\,Hz\,s$^{-1}$ and no finer at any baseline. For the
\TjNOpen{} pairs it cannot separate -- \TjNPowerOpen{} of them resolved and
failed by the constant-acceleration comparison -- nothing measurable tells one
carrier from two.

$\beta$~Pictoris is the closest case, and it is the one the comparison
cannot dismiss. Its two crossings sit at \TjBpNuLo{}
and \TjBpNuHi\,GHz, \TjBpDvKms\,km\,s$^{-1}$ apart in the star's frame, in
blocks \TjBpDtH\,hours apart -- both inside the methanol $5_K\to4_K$ ladder,
which Appendix~\ref{app:linemask} excludes as their source -- with both
fitted drifts negative
($\TjBpDriftLo$ and $\TjBpDriftHi$\,Hz\,s$^{-1}$), so a single carrier would
have had to rise in frequency between the epochs and under constant
acceleration misses by \TjBpNSig{} standard deviations -- but the mean
acceleration it needs, \TjBpAccelReq\,m\,s$^{-2}$, is inside the
\WdAGrid\,m\,s$^{-2}$ ceiling. The pair fixes three numbers, the
acceleration at each epoch and the frequency offset between them, and a
circular orbit about a star of \EcMassBp\,M$_\odot$ \citep{Crifo1997} has
three free parameters: one orbit reproduces all three at once, at
\ObDayBp\,d in period, \ObIncBp$^\circ$ in inclination and \ObAuBp\,au in
radius, eccentric orbits fit up to $e$ of about \ObEccMaxBp{} before the
periastron acceleration leaves the ceiling, and the pair is therefore not
separable.

\begin{table}
%% ★ v4.26 (r25-captions): 5 printed lines -> the author's 3.  Naming
%% Table~\ref{tab:pairs} as what the readings are taken TO says in one clause
%% what two said: "the pairs the trajectory test is asked of" and "the counts
%% in the lower block are those of Table~\ref{tab:pairs}" are the same pairs.
\caption{The \WdNExc{} above-trigger readings of the widened recurrence
window, taken to the pairs of Table~\ref{tab:pairs}; each block of rows is
the one above with one effect removed.}
\label{tab:pairacct}
\footnotesize
%% GENERATED by pairacct_r23.py -- do not hand-edit.
\begin{tabular}{@{}p{0.76\columnwidth}r@{}}
\hline
Covering windows read, over the \PaNCovWinCross{} unattributed crossings a later block covers & \WdNCov \\
\ldots{} holding a statistic above the trigger, in \PaNRepWin{} windows of \PaNRepBlock{} blocks & \WdNExc \\
\hline
Distinct crossing pairs those readings represent & \PaNPairRead \\
\ldots{} readings that are the second reading of such a pair & \PaNPairTwice \\
Pairs separated by more than $\pm\WdHalfKms$\,km\,s$^{-1}$ in the star's frame, to \PaDvBeyondMax\,km\,s$^{-1}$ & \PaNPairBeyond \\
Pairs inside that interval & \PaNPairInside \\
\hline
Pairs with at least one unattributed member & \TjNCandPair \\
\ldots{} with a fitted drift and a block epoch for both members & \TjNPair \\
\hline
\ldots{} separated by the bounded-trajectory criterion & \TjNExcl \\
\ldots{} open & \TjNOpen \\
\hline
\end{tabular}

\end{table}

%% ★★ v4.23: THIS SUBSECTION AND ITS TABLE SIT AT THE END OF THE
%% APPENDIX FOR A REASON THAT IS NOT STYLE.  Placed next to the
%% registration budget, where it belongs by subject, its float took
%% the number Table~D1 and pushed the crossing ledger's sample to D2
%% and the recurrence positive controls to D3 -- and the positive
%% controls are the table Referee 2's first point is about, which he
%% and the response letter both call Table~D2.  Renumbering the
%% table a referee is reading about is a cost with no benefit, so the
%% new float goes last and takes the new number.  \S\ref{sec:edge}
%% is cited from \S4 and \S5.6 and the pointer is unaffected.
\subsection{The drift grid's edges}
\label{sec:edge}

The drift trials of a window run from $-\dot\nu_{\max}$ in equal steps of
half a native channel across the retained track, and the last trial is the
first at or beyond $+\dot\nu_{\max}$, so the interval searched is asymmetric
by up to one step and the ceilings themselves, computed window by window from
Equation~\ref{eq:accel}, reproduce it to $\DriftRelDev$ over all
\DriftRelNWin{} Class~A windows. Each window's grid is recovered from two
quantities that enter no common calculation -- the trial count recorded for
it and the channel width and retained span its own crossing record carries --
and the two agree on \EgNGrid{} of the \EgNDrift{} crossings that carry a
fitted drift, on every one of which the fitted rate is a trial of that grid
exactly; \EgNEdge{} lie on an outermost trial (Table~\ref{tab:edge}), toward
\EgEdgeStars, and are reported as lower bounds. \EgNEdgeExcl{} are among the
\RcNExcl{} crossings \S\ref{sec:etacrv} bounds; the other \EgNEdgeOut{} carry
no amplitude bound, lying in the reserved hold-out, in the quality-failing
block and in the one window retaining no discovery spectrum, which is why
this table reaches above that section's range. The remaining
\EgNNoGrid{} are in archival-tail windows for which no grid is recorded
anywhere.

\begin{table}
%% ★ v4.26 (r25-captions): 6 printed lines -> the author's 3.  DELETED after
%% checking: "a maximum at the boundary of the interval searched is a lower
%% bound" is \S\ref{sec:etacrv}'s own definition of \emph{edge-limited}, and
%% "Table~\ref{tab:pairs} gives them in the star's frame" is now a clause of
%% that same paragraph, which is where a reader meets both tables.
\caption{The crossings whose fitted drift lies on an outermost trial of
their window's own grid, with that grid. Rates are in the observatory
frame.}
\label{tab:edge}
\footnotesize
%% GENERATED by edge_v423.py -- do not hand-edit.
\begin{tabular}{@{}llr@{~}r@{~}c@{~}r@{~}l@{}}
\hline
Star & block & $\nu$ & $\dot\nu_{\rm obs}$ & trial & \multicolumn{2}{c}{$|a_{\rm los}|$} \\
 & & (GHz) & (Hz\,s$^{-1}$) & & \multicolumn{2}{c}{(m\,s$^{-2}$)} \\
\hline
BD+05~1668 & \texttt{X28a8} & 242.195 & $+3452$ & 2/2 & $\ge$ & 4.27 \\
HD~14055 & \texttt{X29ae} & 331.156 & $-3976$ & 0/135 & $\ge$ & 3.60 \\
HD~14055 & \texttt{X2c3bb} & 345.526 & $+4184$ & 141/141 & $\ge$ & 3.63 \\
HD~161868 & \texttt{X10ca8} & 345.634 & $+4195$ & 127/127 & $\ge$ & 3.64 \\
TRAPPIST$-$1 & \texttt{X259b} & 346.245 & $-4623$ & 0/137 & $\ge$ & 4.00 \\
\hline
\end{tabular}

\end{table}
